\section{The Bridge: bindgen Over a Vendor Header}
\label{sec:ble-bridge}

Generating Rust bindings from \codeinline{sl_bt_api.h} is one command and three surprises. All three
of this project's were compile errors in the wrapper, and all three are the kind that look like
mistakes in your own code.

\subsection{Surprise 1: The Union Is Emulated}

The header declares its event payload as an \emph{anonymous} union inside a struct. bindgen
represents anonymous unions with \codeinline{\_\_BindgenUnionField<T>}, regardless of the Rust target
requested — measured: adding \codeinline{.rust\_target(Stable\_1\_71)} left the generated file
byte-identical.

The consequence is 18 errors of the form:

\begin{minted}{text}
error[E0609]: no field `connection` on type
              `&__BindgenUnionField<sl_bt_evt_connection_opened_s>`
\end{minted}

The fix is not a bindgen flag. It is to read the payload the way the emulation requires:

\begin{minted}{rust}
// before: reads the wrapper as if it were the union
let data = &evt.data.evt_connection_opened;

// after: says out loud that the active member is the caller's responsibility
let data = unsafe { evt.data.evt_connection_opened.as_ref() };
\end{minted}

That \codeinline{unsafe} is the correct amount of friction. Reading a union is unsafe because nothing
checks that the member you chose is the one the sender wrote — and the event ID is the only thing that
tells you.

\subsection{Surprise 2: A Flexible Array Member Is Not a Slice}

\begin{minted}{text}
pub struct uint8array {
    pub len: u8,                        // "< Number of bytes stored in data"
    pub data: __IncompleteArrayField<u8>,
}
\end{minted}

\codeinline{data} has no length and no \codeinline{Index}, so \codeinline{data.value.data[i]} does
not compile. The bytes are read through the pointer, with the length taken from the header's own
field:

\begin{minted}{rust}
let byte = unsafe { data.value.data.as_ptr().add(i).read() };
\end{minted}

Two facts worth keeping: the length is authoritative (it was written by the sender), and the pointer
arithmetic is the only way to reach a flexible array — which is exactly why C APIs that return them
need a total-length parameter somewhere.

\subsection{Surprise 3: One Field, Two Names}

\begin{minted}{text}
sl_bt_evt_gatt_server_attribute_value_t      ->  attribute: u16
sl_bt_evt_gatt_server_user_read_request_t    ->  characteristic: u16
\end{minted}

The same 16-bit attribute handle, named differently in two events of the same SDK version. A wrapper
that assumes one name fails to compile on the other, and a wrapper that ``fixes'' it by renaming the
Rust type breaks the one that was right.

The project's rule: convert at the boundary, once, with a comment naming both spellings. The
application-facing type has one field, \codeinline{attribute}, and the FFI layer knows about the
vendor's inconsistency so nothing above it has to.

\subsection{What the Generated Code Is Good For}

With those three handled, the bindings are genuinely useful in the way the SVD's are: every field
carries the header's doc comment, so \codeinline{cargo doc} shows the vendor's own description of the
value you are about to write. That is the strongest argument for generating bindings rather than
hand-writing them — not the function signatures, which are mechanical, but the documentation that
travels with them.
